% =====================================================================
% §3 Problem setting & data
% Adapted from the CS-137 §2 Dataset (light scrub).
% =====================================================================

\section{Problem setting and data}
\label{sec:data}

The task uses three input modalities, all hourly and aligned in UTC: ISO-NE per-zone electricity demand, NOAA HRRR weather rasters, and 44-dimensional calendar one-hot vectors. We describe each in turn, then the train / val / test split rationale and the four-step normalization chain shared across all models. A datasheet-style summary (motivation, composition, collection, preprocessing, uses, distribution, maintenance, PII) is in Appendix~\ref{app:datasheet}.

\subsection{ISO-NE per-zone electricity demand}
\label{sec:data:demand}

ISO New England (ISO-NE) is the regional transmission organisation operating the wholesale electricity market for the six New England states. Hourly system-level and per-zone demand data are published through the public \emph{Energy, Load, and Demand} reports portal \citep{isone2024data}. We use eight zones corresponding to ISO-NE's load-zone definitions:

\begin{itemize}\setlength\itemsep{1pt}
  \item \zone{ME} (Maine), \zone{NH} (New Hampshire), \zone{VT} (Vermont) --- the northern zones, smaller, more rural, more electric-heating-dominant in winter.
  \item \zone{CT} (Connecticut), \zone{RI} (Rhode Island), \zone{SEMA} (south-east Massachusetts), \zone{WCMA} (west-central Massachusetts), \zone{NEMA\_BOST} (north-east Massachusetts / metro Boston) --- the southern / urban-coastal zones, larger, more cooling-driven in summer.
\end{itemize}

Demand is recorded hourly in MWh. Zone-level mean and standard deviation, computed from 500 random training samples drawn from 2019--2020, are listed in Table~\ref{tab:zone_stats}. \zone{VT} is the smallest (mean $\approx$\,577\,MWh/h, std $\approx$\,113), \zone{CT} is the largest (mean $\approx$\,3186\,MWh/h, std $\approx$\,711); a factor-of-5.5 dynamic range across zones. Five calendar years (2019--2023) are used during development; in deployment (Part II) we additionally consume the public 5-minute estimated zonal load feed for the most recent days, rolled up to hourly.

\begin{table}[!h]
\centering
\small
\caption{Per-zone demand statistics (MWh) computed from 500 random training samples, used as the z-score $(\mu_y, \sigma_y)$ for the target during training.}
\label{tab:zone_stats}
\begin{tabular}{@{}lrrrrrrrr@{}}
\toprule
& \zone{ME} & \zone{NH} & \zone{VT} & \zone{CT} & \zone{RI} & \zone{SEMA} & \zone{WCMA} & \zone{NEMA\_BOST} \\
\midrule
$\mu_y$ (MWh) & 1306.7 & 1294.0 & \phantom{0}577.4 & 3185.7 & \phantom{0}879.5 & 1571.1 & 1809.7 & 2661.7 \\
$\sigma_y$ (MWh) & \phantom{0}211.1 & \phantom{0}257.9 & \phantom{0}113.1 & \phantom{0}710.8 & \phantom{0}197.4 & \phantom{0}368.5 & \phantom{0}340.0 & \phantom{0}524.5 \\
\bottomrule
\end{tabular}
\end{table}

\subsection{HRRR weather rasters}
\label{sec:data:weather}

NOAA's High-Resolution Rapid Refresh (HRRR) is the operational North-American mesoscale numerical-weather-prediction system on a 3\,km grid \citep{benjamin2016hrrr,noaa2024hrrr}. We use the f00 analysis (instantaneous state at the valid hour), restricted to a New England bounding box ($40.5{-}47.5\,^{\circ}$N, $-74.0$ to $-66.0\,^{\circ}$E) and bilinearly regridded from HRRR's native Lambert conformal grid to a regular $450\times 449$ lat/lon grid (resolution $\sim$1.6\,km/cell). At training time we use 24 separate $f00$ analyses for both the history and future windows (this implies privileged access to ground-truth future weather; we discuss the implication for live deployment in \S\ref{sec:deploy:arch}). Seven channels per timestep, in fixed order, with z-score statistics from 500 random training samples:

\begin{table}[!h]
\centering
\small
\caption{HRRR weather channels and per-channel z-score statistics from 500 random training samples. \emph{Caveat on \texttt{DSWRF} and \texttt{APCP\_1hr}:} the cached means below (0.09 and 158.13) reflect the 500-sample draw on the cached HRRR rasters; the absolute magnitudes do not match a typical New England annual climatology in the unit labels shown, indicating an upstream unit-or-accumulation-period idiosyncrasy in the HRRR fetcher (see Appendix~\ref{app:hrrr}). Because the model is trained against these specific cached statistics, the four-step normalization chain (\S\ref{sec:data:norm}) is internally consistent regardless of the external unit interpretation.}
\label{tab:weather_stats}
\begin{tabular}{@{}llrr@{}}
\toprule
Channel & Description & $\mu_x$ & $\sigma_x$ \\
\midrule
\texttt{TMP}      & 2-meter air temperature (K)                         & 281.17 & \phantom{00}11.31 \\
\texttt{RH}       & 2-meter relative humidity (\%)                      & \phantom{0}74.42 & \phantom{00}16.05 \\
\texttt{UGRD}     & 10-meter U wind component (m/s)                     & \phantom{00}7.76 & \phantom{000}4.55 \\
\texttt{VGRD}     & 10-meter V wind component (m/s)                     & \phantom{00}1.45 & \phantom{000}3.84 \\
\texttt{GUST}     & surface wind gust (m/s)                             & \phantom{00}0.42 & \phantom{000}3.73 \\
\texttt{DSWRF}    & downward shortwave radiation flux (W/m\textsuperscript{2}) & \phantom{00}0.09 & \phantom{000}0.55 \\
\texttt{APCP\_1hr}& 1-hour accumulated precipitation (kg/m\textsuperscript{2})  & 158.13 & \phantom{0}251.87 \\
\bottomrule
\end{tabular}
\end{table}

Each timestep produces a tensor of shape $(450, 449, 7)$ stored as a single \texttt{.pt} file ($\sim$5.7\,MB float32 / $\sim$2.9\,MB float16). The full 2019--2023 archive is $\sim$130\,GB compressed, fetched from NOAA's S3 mirror \texttt{s3://noaa-hrrr-bdp-pds/} with the \texttt{herbie-data} library \citep{blaylock2023herbie}; the fetcher and regrid procedure are documented in Appendix~\ref{app:hrrr} and re-used unchanged for live deployment.

\subsection{Calendar features}
\label{sec:data:cal}

For every hour we build a 44-dimensional one-hot calendar vector $C_t\in\{0,1\}^{44}$:
\begin{itemize}\setlength\itemsep{1pt}
  \item 24 dimensions for hour-of-day (0--23, one-hot).
  \item \phantom{0}7 dimensions for day-of-week (Monday--Sunday, one-hot).
  \item 12 dimensions for month-of-year (Jan--Dec, one-hot).
  \item \phantom{0}1 dimension for US-holiday flag (NYSE schedule).
\end{itemize}
These features capture human-activity periodicity --- commuting peaks, weekend troughs, holiday dips --- which demand profiles strongly track and which weather alone does not encode.

\subsection{Train / validation / test splits}
\label{sec:data:splits}

Table~\ref{tab:splits} summarises the splits used during development. The CNN-Transformer baseline (\S\ref{sec:baseline}) trains on 2019--2020 and validates on 2021; the encoder-decoder variants (\S\ref{sec:archsearch}) train on 2019--2021 and validate on 2022. All training-time MAPE numbers reported in this paper are computed on the last two days of 2022.

\begin{table}[!h]
\centering
\small
\caption{Train / validation / test splits used in development.}
\label{tab:splits}
\begin{tabular}{@{}llr@{}}
\toprule
Split & Years & Approx.\ samples \\
\midrule
Train (baseline)             & 2019--2020              & $\sim$17\,K \\
Train (encoder-decoder)      & 2019--2021              & $\sim$26\,K \\
Validation (baseline)        & 2021                    & $\sim$8\,K  \\
Validation (encoder-decoder) & 2022                    & $\sim$9\,K  \\
Self-eval test               & last 2 days of 2022     & 2 windows   \\
Live deployment              & rolling last 7 days (2026+) & 7 windows / day \\
\bottomrule
\end{tabular}
\end{table}

\subsection{Four-step normalization pipeline}
\label{sec:data:norm}

All models share the same four-step normalization chain so training-time predictions and live deployment-time predictions live in the same reference frame:

\begin{enumerate}\setlength\itemsep{1pt}
  \item \textbf{Compute statistics once} (\texttt{compute\_norm\_stats()} in the dataset module): pull 500 random training samples, compute per-channel $(\mu_x, \sigma_x)$ for the 7 weather channels and per-zone $(\mu_y, \sigma_y)$ for the 8 demand zones. Cache to \texttt{norm\_stats.pt} alongside the checkpoint.
  \item \textbf{Z-score at training and inference time}: weather inputs are normalized per-channel via $X' = (X - \mu_x) / \sigma_x$; demand inputs and targets are normalized per-zone via $Y' = (Y - \mu_y) / \sigma_y$. The model trains and predicts entirely in z-score space.
  \item \textbf{Train against MSE} in z-score space (so all 8 zones contribute on equal footing despite the factor-of-5.5 dynamic range).
  \item \textbf{De-normalize for evaluation}: predictions are mapped back to physical MWh via $\hat Y = \hat Y' \cdot \sigma_y + \mu_y$ before MAPE is computed (Equation~\ref{eq:mape}).
\end{enumerate}

The same checkpoint-embedded \texttt{norm\_stats.pt} is loaded by the live serving stack (\S\ref{sec:deploy:arch}), so a forward pass on real-time inputs produces predictions on the same MWh scale as the training-time evaluation.
